% GENERATED FILE. Edit canonical lessons, then run npm run generate:books.
% canonical-source-hash: fnv1a64:fe76828788ec1fab
% canonical-lessons: AR-C471-khalij, AR-C471-muhit, AR-C471-marjan, AR-C471-tayyar, AR-C471-umq

\chapter{Bays and oceans}
\label{ch:ar-bays-and-oceans}
\hlchaptermodality{471}

\begin{chapteropening}
\textbf{By the end of this chapter:} \emph{I can say a gulf, an ocean, coral, a current and depth.}

Complete the last lesson of chapter 471: I can say a gulf, an ocean, coral, a current and depth.
\end{chapteropening}

\section[\texorpdfstring{\ar{خليج} (khalīj)}{خليج (khalīj)}]{\ar{خليج} (khalīj) — a gulf; a bay}

\label{lesson:AR-C471-khalij}



\begin{quote}
Before the new one: say the Arabic for a riverbank, then the Arabic for sweet.
\end{quote}

\subsection*{You'll want to know: \ar{خليج}}
\textbf{\ar{خليج}} — \emph{khalīj} — \textquotedblleft{}a gulf; a bay\textquotedblright{}.

\begin{quote}
\ar{سبحنا} \ar{في} \ar{خليج} \ar{صغير}.
\end{quote}

\begin{quote}
\emph{sabaḥnā fī khalīj ṣaghīr.}
\end{quote}

\begin{quote}
We swam in a small bay.
\end{quote}

\begin{cousinweb}
\textbf{\ar{خليج}} (\emph{khalīj}) is built on the native Arabic root \emph{kh-l-j}, \textquotedblleft{}to draw away\textquotedblright{} (Lane). A \emph{khalīj} was a channel drawn off a river, and then an arm of the sea.
\end{cousinweb}

\begin{grammarlens}[title={What you've built}]
One more word for getting around.
\end{grammarlens}

\subsection*{Your turn}
\begin{itemize}
\raggedright
  \item \emph{Say it:} \emph{khalīj}
  \item \emph{Say it:} \emph{khalīj}, once more
  \item \emph{Say it:} \emph{sabaḥnā fī khalīj ṣaghīr.} — We swam in a small bay.
  \item \emph{Recall:} say the Arabic for a riverbank, then the Arabic for sweet, then say \emph{khalīj} again
\end{itemize}

\begin{tcolorbox}[breakable,colback=teal!4,colframe=teal!35!black,title={Before you move on}]
What does \ar{خليج} mean? (\textquotedblleft{}A gulf; a bay\textquotedblright{}.) Say it once more.
\end{tcolorbox}
\section[\texorpdfstring{\ar{محيط} (muḥīṭ)}{محيط (muḥīṭ)}]{\ar{محيط} (muḥīṭ) — an ocean}

\label{lesson:AR-C471-muhit}



\begin{quote}
Before the new one: say the Arabic for sweet, then the Arabic for a gulf.
\end{quote}

\subsection*{You'll want to know: \ar{محيط}}
\textbf{\ar{محيط}} — \emph{muḥīṭ} — \textquotedblleft{}an ocean\textquotedblright{}.

\begin{quote}
\ar{عبرت} \ar{السفينة} \ar{المحيط} \ar{في} \ar{شهر}.
\end{quote}

\begin{quote}
\emph{ʿabarat as-safīna al-muḥīṭ fī shahr.}
\end{quote}

\begin{quote}
The ship crossed the ocean in a month.
\end{quote}

\begin{cousinweb}
\textbf{\ar{محيط}} (\emph{muḥīṭ}) is built on the native Arabic root \emph{ḥ-w-ṭ}, \textquotedblleft{}to surround\textquotedblright{} (Lane). It is the Form IV participle: the sea that surrounds the land.
\end{cousinweb}

\begin{grammarlens}[title={What you've built}]
One more word for getting around.
\end{grammarlens}

\subsection*{Your turn}
\begin{itemize}
\raggedright
  \item \emph{Say it:} \emph{muḥīṭ}
  \item \emph{Say it:} \emph{muḥīṭ}, once more
  \item \emph{Say it:} \emph{ʿabarat as-safīna al-muḥīṭ fī shahr.} — The ship crossed the ocean in a month.
  \item \emph{Recall:} say the Arabic for sweet, then the Arabic for a gulf, then say \emph{muḥīṭ} again
\end{itemize}

\begin{tcolorbox}[breakable,colback=teal!4,colframe=teal!35!black,title={Before you move on}]
What does \ar{محيط} mean? (\textquotedblleft{}An ocean\textquotedblright{}.) Say it once more.
\end{tcolorbox}
\section[\texorpdfstring{\ar{مرجان} (marjān)}{مرجان (marjān)}]{\ar{مرجان} (marjān) — coral}

\label{lesson:AR-C471-marjan}



\begin{quote}
Before the new one: say the Arabic for a gulf, then the Arabic for an ocean.
\end{quote}

\subsection*{You'll want to know: \ar{مرجان}}
\textbf{\ar{مرجان}} — \emph{marjān} — \textquotedblleft{}coral\textquotedblright{}.

\begin{quote}
\ar{رأينا} \ar{المرجان} \ar{تحت} \ar{الماء}.
\end{quote}

\begin{quote}
\emph{raʾaynā al-marjān taḥta al-māʾ.}
\end{quote}

\begin{quote}
We saw the coral under the water.
\end{quote}

\begin{cousinweb}
\textbf{\ar{مرجان}} (\emph{marjān}) was borrowed from Syriac (EtymDB).
\end{cousinweb}

\begin{grammarlens}[title={What you've built}]
One more word for getting around.
\end{grammarlens}

\subsection*{Your turn}
\begin{itemize}
\raggedright
  \item \emph{Say it:} \emph{marjān}
  \item \emph{Say it:} \emph{marjān}, once more
  \item \emph{Say it:} \emph{raʾaynā al-marjān taḥta al-māʾ.} — We saw the coral under the water.
  \item \emph{Recall:} say the Arabic for a gulf, then the Arabic for an ocean, then say \emph{marjān} again
\end{itemize}

\begin{tcolorbox}[breakable,colback=teal!4,colframe=teal!35!black,title={Before you move on}]
What does \ar{مرجان} mean? (\textquotedblleft{}Coral\textquotedblright{}.) Say it once more.
\end{tcolorbox}
\section[\texorpdfstring{\ar{تيار} (tayyār)}{تيار (tayyār)}]{\ar{تيار} (tayyār) — a current (of water or air)}

\label{lesson:AR-C471-tayyar}



\begin{quote}
Before the new one: say the Arabic for an ocean, then the Arabic for coral.
\end{quote}

\subsection*{You'll want to know: \ar{تيار}}
\textbf{\ar{تيار}} — \emph{tayyār} — \textquotedblleft{}a current (of water or air)\textquotedblright{}.

\begin{quote}
\ar{كان} \ar{التيار} \ar{في} \ar{النهر} \ar{قويا}.
\end{quote}

\begin{quote}
\emph{kāna at-tayyār fī an-nahr qawiyyan.}
\end{quote}

\begin{quote}
The current in the river was strong.
\end{quote}

An electric current is \emph{tayyār kahrabāʾī}.

\begin{cousinweb}
\textbf{\ar{تيار}} (\emph{tayyār}) is listed by Lane under \emph{t-y-r} as the waves of a sea or a great river; the origin is not settled.
\end{cousinweb}

\begin{grammarlens}[title={What you've built}]
One more word for getting around.
\end{grammarlens}

\subsection*{Your turn}
\begin{itemize}
\raggedright
  \item \emph{Say it:} \emph{tayyār}
  \item \emph{Say it:} \emph{tayyār}, once more
  \item \emph{Say it:} \emph{kāna at-tayyār fī an-nahr qawiyyan.} — The current in the river was strong.
  \item \emph{Recall:} say the Arabic for an ocean, then the Arabic for coral, then say \emph{tayyār} again
\end{itemize}

\begin{tcolorbox}[breakable,colback=teal!4,colframe=teal!35!black,title={Before you move on}]
What does \ar{تيار} mean? (\textquotedblleft{}A current (of water or air)\textquotedblright{}.) Say it once more.
\end{tcolorbox}
\section[\texorpdfstring{\ar{عمق} (ʿumq)}{عمق (ʿumq)}]{\ar{عمق} (ʿumq) — depth}

\label{lesson:AR-C471-umq}



\begin{quote}
Before the new one: say the Arabic for coral, then the Arabic for a current.
\end{quote}

\subsection*{You'll want to know: \ar{عمق}}
\textbf{\ar{عمق}} — \emph{ʿumq} — \textquotedblleft{}depth\textquotedblright{}.

\begin{quote}
\ar{عمق} \ar{البحيرة} \ar{كبير}.
\end{quote}

\begin{quote}
\emph{ʿumq al-buḥayra kabīr.}
\end{quote}

\begin{quote}
The depth of the lake is great.
\end{quote}

\begin{cousinweb}
\textbf{\ar{عمق}} (\emph{ʿumq}) is built on the native Arabic root \emph{ʿ-m-q}, \textquotedblleft{}to be deep\textquotedblright{} (Lane). It is the noun beside \emph{ʿamīq}, \textquotedblleft{}deep\textquotedblright{}.
\end{cousinweb}

\begin{grammarlens}[title={What you've built}]
One more word for getting around.
\end{grammarlens}

\subsection*{Your turn}
\begin{itemize}
\raggedright
  \item \emph{Say it:} \emph{ʿumq}
  \item \emph{Say it:} \emph{ʿumq}, once more
  \item \emph{Say it:} \emph{ʿumq al-buḥayra kabīr.} — The depth of the lake is great.
  \item \emph{Recall:} say the Arabic for coral, then the Arabic for a current, then say \emph{ʿumq} again
\end{itemize}

\begin{tcolorbox}[breakable,colback=teal!4,colframe=teal!35!black,title={Before you move on}]
What does \ar{عمق} mean? (\textquotedblleft{}Depth\textquotedblright{}.) Say it once more.
\end{tcolorbox}
